
\section{The ScaleMaster Dataset}
\label{sec:scalemaster_dataset}

To address the inadequacy of existing benchmarks for evaluating robustness for scale consistency in monocular visual \ac{SLAM}, we introduce and publicly release a new challenging dataset, \textbf{ScaleMaster Dataset}. The baseline metric SE(3) poses for this dataset were acquired using Apple’s ARKit framework, an off-the-shelf state-of-the-art method known to operate robustly with centimeter-level accuracy even in large-scale indoor environments spanning several hundred meters. We additionally performed manual inspections to ensure that the reconstructed maps were well aligned. The detailed sequence information is summarized in \tabref{dataset_summary}, while \tabref{dataset_comparison} provides a comparison against standard benchmark datasets that have been widely used for evaluating dense visual SLAM systems. For brevity, we present a representative subset of our 25 sequences in \tabref{dataset_summary}. The complete details for all sequences are provided in our supplementary material.

A key contribution of our benchmark is the provision of dense, high-fidelity 3D ground truth maps. We constructed the 3D reference maps by projecting a high-precision 3D LiDAR point cloud onto these ARKit trajectories.

% We acknowledge that while the iPhone ARKit trajectory is known to operate robustly with centimeter-level accuracy for local tracking, it is not entirely free from long-term drift. Therefore, we do not claim our ground truth maps possess metrology-grade, global accuracy. However, our primary goal is to evaluate the scale consistency and geometric distortion of monocular SLAM reconstructions, rather than their absolute global pose accuracy. For this purpose, the smooth and locally precise trajectory provided by ARKit serves as a strong foundation for generating a reference map.
While ARKit trajectory may still suffer from long-term drift, it provides smooth and locally precise trajectories. As shown in \figref{fig:system_overview}, projecting LiDAR points onto ARKit poses yields maps that are qualitatively well aligned, making them a reliable reference for evaluating scale consistency.

The ScaleMaster dataset comprises 25 sequences and is intentionally designed to include a variety of challenging scenarios rarely found in other datasets. The dataset is structured to support different evaluation depths: an initial set of 18 sequences is provided with trajectory-only ground truth, while a core subset of 7 sequences features both trajectories and the dense, high-fidelity 3D ground truth maps. These include multi-floor structures with significant vertical motion, long trajectories with repetitive views, and large spaces that exceed the effective sensing range of typical RGB-D sensors.

For ground truth data acquisition, we utilized a custom-designed handheld rig (see the 'Hardware' panel in \figref{fig:system_overview}) equipped with an iPhone 14 Pro, a HAP LiDAR sensor, and a Gemini335L camera, which ensures precise temporal synchronization between the captured imagery and LiDAR scans, thereby improving data quality and reliability.

These characteristics establish our dataset as a novel stress-test benchmark for systematically evaluating the true robustness of modern SLAM algorithms, with a particular focus on \text{scale-related challenges} such as ambiguity and inconsistency that emerge in diverse, real-world environments.
